% !TeX root = ../PIFB2.tex
\section{Derivational and Source Methods}\label{sec:methods}
\label{sec:methods_metagent}
The manuscript uses analytic probability identities, typed geometric constructions and explicit finite-dimensional examples. New algebraic steps were checked independently, with symbolic calculations used for finite identities rather than as substitutes for general proofs. The formation-selection result is cited from its complete companion derivation with its fixed-temperature and almost-everywhere hypotheses preserved. No new empirical experiment or physical validation is reported.

\subsection{Source use}
Primary articles and author manuscripts were consulted for the information-geometric, filtering, coarse-graining and Lorentz constructions. Literature-assisted research was used to identify candidate extensions and contrary cases. A suggested citation was not treated as proof that its hypotheses apply to the present model. The source inventory distinguishes adopted references, inherited citations and unresolved access limitations.

\subsection{Language-model assistance}
Language models assisted with literature discovery, mathematical reconstruction, symbolic checking and manuscript organization and drafting under author direction. Exa Agent was used at ultra effort for the three focused research questions on inference and formation, Lorentz structure, and persistence and inertia. Claims retained in the manuscript rest on their stated derivations and primary sources, not on agreement among models.

